\emph{(i)} Substitute $u = u_{0} + \lambda x$, $v = u_{0} + \lambda y$: $\int_{I}\int_{I}\lambda^{4p}|R(u_{0} + \lambda x, u_{0} + \lambda y)|^{2}\,dx\,dy = \lambda^{4p - 2}\int_{u_{0} + \lambda I}\int_{u_{0} + \lambda I}|R(u, v)|^{2}\,du\,dv$. The factor is $1$ for all $\lambda$ exactly when $p = \tfrac12$; this is the weight of the boost on a generator and the only weight preserving the anticommutator (Theorem \claimref{synthesis.stationary-generator-states-and-the-edge-scaling}(i)). The check confirms the identity with exact integrals of polynomial kernels over rational rectangles at $\lambda = t^{2}$ and $p = k/4$, equality for $k = 2$ and the factor $t^{2k}/\lambda^{2}$ otherwise, and that the boost rotor of $\Cl(1,3)$ dilates the generator about its edge.

\emph{(ii)} By Fact \claimref{synthesis.imported-powers-stormer-takesaki-disjointness}(a), applied to the CAR algebra over $L^{2}(J)\otimes K$ with the compressed symbols $A$, $B$ ($0 \le A, B \le 1$, kernels the restrictions of $W_{\omega}$ and $W_{0}$, up to the convention of which of $S$ and $1 - S$ carries the two-point function, which does not affect the differences), local normality on $J$ makes $A^{1/2} - B^{1/2}$ Hilbert-Schmidt; then $A - B = A^{1/2}(A^{1/2} - B^{1/2}) + (A^{1/2} - B^{1/2})B^{1/2}$ with $\|A^{1/2}\|, \|B^{1/2}\| \le 1$ gives $\|A - B\|_{\mathrm{HS}} \le 2\|A^{1/2} - B^{1/2}\|_{\mathrm{HS}}$ (the identity and the inequality are checked exactly on rational contractions), so $R$ is square integrable on $J \times J$. Now $W_{0}$ is invariant under translations and under the weight-$\tfrac12$ dilations, so $\lambda W_{\omega}(u_{0} + \lambda x, u_{0} + \lambda y) - W_{0}(x, y) = (S_{\lambda}R)(x, y)$ with $p = \tfrac12$. For a bounded interval $I$ and $\lambda$ small, $u_{0} + \lambda I \subset J$, and by (i) $\|S_{\lambda}R\|_{L^{2}(I \times I)} = \|R\|_{L^{2}((u_{0} + \lambda I)^{2})}$, which tends to $0$ by dominated convergence, the sets shrinking to the point $(u_{0}, u_{0})$. Hilbert-Schmidt convergence implies operator-norm convergence of the compressed symbols, and the $n$-point functions of a gauge-invariant quasi-free state are determinants of matrix elements of its symbol, so they converge on test functions supported in $I$. Hence the scaling limit at $u_{0}$ exists and is $\omega_{0}$, in the sense of Theorem \claimref{synthesis.stationary-generator-states-and-the-edge-scaling}. For a bounded remainder the shrunk norm is $O(\lambda)$ (checked: $\le \lambda^{2}|I|^{2}\max|R|^{2}$ for the squared norm on polynomial samples), which is the case of Theorem \claimref{synthesis.null-vacuum-singularity-under-reparametrisation-and-dilation}(ii).

\emph{(iii)} By Theorem \claimref{synthesis.null-vacuum-singularity-under-reparametrisation-and-dilation}(i), the vacuum kernels of two regular parameters differ, at weight $\tfrac12$, by a kernel continuous across the diagonal, hence bounded and square integrable on bounded intervals, and the weight-$\tfrac12$ change of variables is unitary on $L^{2}$; so the local Hilbert-Schmidt condition relative to one regular vacuum is equivalent to it relative to any other. The exact example: with $f' = (1 + u^{2})^{2}$, $(1 + u^{2})(1 + v^{2}) - q(u, v) = (u - v)^{2}\big(\tfrac13 - \tfrac15(u^{2} + 3uv + v^{2})\big)$, checked on a rational grid; the leading coefficient $\tfrac13$ is $S_{f}(0)/12$ with $S_{f}(0) = 4$ computed from the polynomial derivatives.

\emph{(iv)} Stationarity and the dilation form of the boost give $R(u, v) = sR(su, sv)$ for all $s > 0$. Let $R$ be square integrable on $(-a, a)^{2}$. By (i) with $u_{0} = 0$, $\|R\|_{L^{2}((-a, a)^{2})} = \|R\|_{L^{2}((-sa, sa)^{2})}$ for every $s \in (0, 1]$, and the right side tends to $0$ as $s \to 0$ by dominated convergence. So $R = 0$ almost everywhere on $(-a, a)^{2}$, and by the homogeneity, which maps any point into $(-a, a)^{2}$, everywhere. A gauge-invariant quasi-free state is fixed by its two-point function, so $\omega = \omega_{0}$, which satisfies (N$_{+}$) (Theorem \claimref{synthesis.null-vacuum-on-a-horizon-generator}(ii)); in its GNS representation the affine translations have nonnegative generator.

\emph{(v)} If $R$ is homogeneous of degree $-1$, $|R|^{2}$ is homogeneous of degree $-2$, and the substitution $(u, v) \mapsto (2u, 2v)$ maps the shell $\{2^{-k-1}a \le \max(|u|, |v|) < 2^{-k}a\}$ onto the next with the same integral of $|R|^{2}$; if $R$ is nonzero on a set of positive measure, some shell, hence every shell, has a positive integral, and the sum over shells diverges. For the cross-edge block $\tfrac{I_{4}}{2\pi}(u - v)^{-1}$, $u > 0 > v$, $w = -v$: $\int_{a}^{b}\int_{c}^{d}(u + w)^{-2}\,dw\,du = \ln\frac{(a + d)(b + c)}{(a + c)(b + d)}$, and on the squares $[s, 2s]^{2}$ the argument is $\tfrac98$ for every $s$ (checked exactly, with the multiplicativity of the arguments over adjacent rectangles, which is the additivity of the integral). The cross-edge block of every $\omega_{c}^{\times}$ is $0$, so its difference from $\omega_{0}$ there is minus this block, whatever $c$; on a half at $c = \tfrac12$, with $u = a^{2}$, $v = b^{2}$, the difference is $-\tfrac{I_{4}}{2\pi}(a - b)/((a + b)(a^{2} + b^{2}))$, homogeneous of degree $-1$ in $(u, v)$ and nonzero off the diagonal (checked). More generally, for any stationary $\omega \neq \omega_{0}$ the difference $R$ is homogeneous of degree $-1$ by stationarity and nonzero, so by (iv) it is not square integrable near the edge; by (ii) the restrictions to the algebra of any interval containing the edge are not quasi-equivalent, and by Fact \claimref{synthesis.imported-powers-stormer-takesaki-disjointness}(a) (gauge-invariant quasi-free states are factor states) they are disjoint.

\emph{Reading.} The equivalence within the stationary class is (iv) with the trivial converse ($\omega_{0}$ is locally normal to itself). That a state KMS at temperature $T$ for the affine translations has a remainder bounded on bounded sets follows from its kernel $\tfrac{I_{4}}{2\pi}\pi T/\sinh(\pi T(u - v + I_{4}0))$, whose difference from $W_{0}$ is analytic across the diagonal (argued; it is not in the check), so the local Hilbert-Schmidt condition holds at every point and (ii) applies; local normality of that state is not derived here. The equivalence is stated for the generator-algebra analogue of local normality; that the vacua of two regular parameters are locally quasi-equivalent, not only within a bounded kernel of each other, would need the implementability of the reparametrisation (known for the free chiral fermion, cited from memory) and is not used. The principle whose generator-algebra analogue is adopted, local normality of physical states with respect to the vacuum for the algebras of spacetime regions, and its antecedent, the Hawking temperature from a scaling-limit condition at the horizon, are Fact \claimref{synthesis.imported-haag-local-normality} (Haag; Haag, Narnhofer and Stein 1984); the transfer to the algebra of an interval of one generator is the analogue, not part of that Fact. Literature from memory: K. Fredenhagen and R. Haag, Commun. Math. Phys. 127 (1990) 273.
